% 定位：sec:numerical-results 的“成本对照与结构总结”小节。
% 替换静态说明；保留成本表和 NumOverallFinding 的一次调用。
表\ref{tab:numerical-scenarios}中的 $J_{\mathrm{ref}}$
为解析公式的数值评价，$J_{\mathrm{OCL}}$
为数值优化返回控制按式\eqref{eq:numerical-cost-sum}计算的成本。
成本接近、阶段顺序一致和切换位置接近是相互补充的比较指标，
其中任一项都不能单独替代其余指标。
数值一致性不要求孤立切换瞬间的控制代表值相同。

数值控制允许存在与网格尺度相关的过渡单元，
故切换误差应与相应网格宽度一并解释。
配点状态与重积分状态之间的偏差反映动力学离散的一致性，
并不等同于数值状态与解析最优轨道之间的误差。
若离散成本略低于解析成本，应同时检查原控制的重积分容量约束，
不能将存在容量超出的离散结果称为优于解析最优策略的严格可行控制。

% 只允许实际结果驱动；不得从预期结构拼出“识别结构”。
\NumFindingParagraph{\NumOverallFinding}
